
   \begin{longtable}[c]{|p{1.5cm}|p{1.9cm}|p{1.5cm}|p{1.5cm}|p{1.5cm}|p{1.5cm}|p{1.0cm}|p{1.0cm}|p{1.0cm}| }
   \caption{IO MUX 管脚功能}
    \label{tab:iomuxgpio-pin-functions} \\
        \hline
        \rowcolor{lightgray}
\textbf{管脚编号} & \textbf{管脚名称} & \textbf{功能 0} & \textbf{功能 1} & \textbf{功能 2} & \textbf{功能 3} & \textbf{驱动强度} & \textbf{复位} & \textbf{说明} \\
 \hline
\endfirsthead
        \hline
        \rowcolor{lightgray}
        \textbf{管脚编号} & \textbf{管脚名称} & \textbf{功能 0} & \textbf{功能 1} & \textbf{功能 2} & \textbf{功能 3} & \textbf{驱动强度} & \textbf{复位} & \textbf{说明} \\
        \hline
        \endhead
        \endfoot
4 & XTAL\_32K\_P     & GPIO0         & GPIO0   & -          & -            & 2           & 0 & R \\ \hline
5 & XTAL\_32K\_N     & GPIO1         & GPIO1   & -          & -            & 2           & 0 & R \\ \hline
6 & GPIO2            & GPIO2         & GPIO2   & FSPIQ      & -            & 2           & 1 & R \\ \hline
8 & GPIO3            & GPIO3         & GPIO3   & -          & -            & 2           & 1 & R \\ \hline
9 & MTMS             & MTMS          & GPIO4   & FSPIHD     & -            & 2           & 1 & R \\ \hline
10 & MTDI             & MTDI          & GPIO5   & FSPIWP     & -            & 2           & 1 & R \\ \hline
12 & MTCK             & MTCK          & GPIO6   & FSPICLK    & -            & 2           & 1* & G \\ \hline
13 & MTDO             & MTDO          & GPIO7   & FSPID      & -            & 2           & 1 & G \\ \hline
14 & GPIO8            & GPIO8         & GPIO8   & -          & -            & 2           & 1 & - \\ \hline
15 & GPIO9            & GPIO9         & GPIO9   & -          & -            & 2           & 3 & - \\ \hline
16 & GPIO10          &  GPIO10       & GPIO10  & FSPICS0    & -            & 2           & 1 & G \\ \hline
18 & VDD\_SPI        &  GPIO11       & GPIO11  & -          & -            & 2           & 0 & - \\ \hline
19 & SPIHD           &  SPIHD        & GPIO12  & -          & -            & 2           & 3 & - \\ \hline
20 & SPIWP           &  SPIWP        & GPIO13  & -          & -            & 2           & 3 & - \\ \hline
21 & SPICS0          &  SPICS0       & GPIO14  & -          & -            & 2           & 3 & - \\ \hline
22 & SPICLK          &  SPICLK       & GPIO15  & -          & -            & 2           & 3 & - \\ \hline
23 & SPID            &  SPID         & GPIO16  & -          & -            & 2           & 3 & - \\ \hline
24 & SPIQ            &  SPIQ         & GPIO17  & -          & -            & 2           & 3 & - \\ \hline
25 & GPIO18          &  GPIO18       & GPIO18  & -          & -            & 3           & 0 & USB, G \\ \hline
26 & GPIO19          &  GPIO19       & GPIO19  & -          & -            & 3           & 0* & USB \\ \hline
27 & U0RXD           &  U0RXD        & GPIO20  & -          & -            & 2           & 3 & G \\ \hline
28 & U0TXD           &  U0TXD        & GPIO21  & -          & -            & 2           & 4 & - \\ \hline

    \end{longtable}
\normalsize

\textbf{驱动强度}

``驱动强度" 一栏所示为每个管脚复位后的默认驱动强度。
\begin{itemize}
    \item GPIO2, GPIO3, GPIO4, GPIO5, GPIO18, GPIO19
    \begin{itemize}
        \item \textbf{0} - 驱动电流 = \verb+~+5 mA
        \item \textbf{1} - 驱动电流 = \verb+~+20 mA
        \item \textbf{2} - 驱动电流 = \verb+~+10 mA
        \item \textbf{3} - 驱动电流 = \verb+~+40 mA
    \end{itemize}
    \item 其他 GPIO
    \begin{itemize}
        \item \textbf{0} - 驱动电流 = \verb+~+5 mA
        \item \textbf{1} - 驱动电流 = \verb+~+10 mA
        \item \textbf{2} - 驱动电流 = \verb+~+20 mA
        \item \textbf{3} - 驱动电流 = \verb+~+40 mA
    \end{itemize}
\end{itemize}

\textbf{复位}

``复位" 一栏所示为每个管脚复位后的默认配置。

\begin{itemize}
\item \textbf{0} - IE = 0（输入关闭）
\item \textbf{1} - IE = 1（输入使能）
\item \textbf{2} - IE = 1，WPD = 1（输入使能，下拉电阻使能）
\item \textbf{3} - IE = 1，WPU = 1 （输入使能，上拉电阻使能）
\item \textbf{4} - OE = 1, WPU = 1（输出使能，上拉电阻使能）
\item \textbf{0*} - IE = 0，WPU = 0，GPIO19 的 USB 上拉默认值为 1，因此，其上拉电阻使能，具体见\textbf{说明}。
\item \textbf{1*} - 如果 EFUSE\_DIS\_PAD\_JTAG = 1，则 MTCK 管脚复位后浮空，即 IE = 1。如果 EFUSE\_DIS\_PAD\_JTAG = 0，则 MTCK 管脚连接内部上拉电阻，即 IE = 1，WPU = 1。  
\end{itemize}

\textbf{说明}

\begin{itemize}
\item \textbf{R} - 代表位于 VDD3P3\_RTC 电源域的管脚，部分具有模拟功能，见表 \ref{tab:iomuxgpio-iomax-pins-analog-functions}。
\item \textbf{USB} - GPIO18、GPIO19 为 USB 管脚。USB 管脚的上拉控制由管脚上拉和 USB 上拉共同控制。当其中任意一个为 1 时，对应管脚上拉电阻使能。USB 上拉值对应寄存器 USB\_SERIAL\_JTAG\_DP\_PULLUP。
\item \textbf{G} - 管脚在芯片上电过程中有毛刺，具体见表 \ref{tab:glitch}。

\end{itemize}

\begin{longtable}[c]{ | l | L{3cm} | r |  }
    \caption{芯片上电过程中的管脚毛刺}
    \label{tab:glitch}\\\hline
    \rowcolor{lightgray}
                                           && \multicolumn{1}{c|}{\textbf{典型持续时间}} \\
    \rowcolor{lightgray}
    \multirow{-2}{*}{\textbf{管脚}}    &
    \multirow{-2}{*}{\textbf{毛刺类型}}  & \multicolumn{1}{c|}{\textbf{(ns)}} \\\hline
    MTCK   & 低电平毛刺 & 5     \\ \hline
    MTDO   & 低电平毛刺 & 5     \\ \hline
    GPIO10 & 低电平毛刺 & 5     \\ \hline
    U0RXD  & 低电平毛刺 & 5     \\ \hline
    GPIO18 & 高电平毛刺    & 50000 \\ \hline 

\end{longtable}
